\section{Problem and Information Structure}
\label{sec:problem}
Consider a bidirectional bus service with a fixed fleet and a fixed set of trips $\mathcal T$. Trip $j$ has a public reference departure $p_j$, an operationally published release $\tau_j$, and an actual departure $d_j$. Public references remain unchanged when the controller revises a release. Each controllable stop visit $e$ receives a holding command $h_e$ from a finite legal set $\mathcal H_e$. Vehicle returns determine the supply available to later trips, so the feasible schedules and holding sequences are coupled.

Let $W$ and $R$ denote total passenger waiting and on-board time, respectively, in person-seconds. The accounting includes the prescribed tail of service and uses the same generated passenger population for paired policies. Let $D_{\mathrm{pub}}$ be total positive departure delay measured against the immutable public references. The implemented upper objective is
\begin{equation}
 J_0=\frac{W+R}{3{,}600{,}000}+\frac{D_{\mathrm{pub}}}{94{,}320}.
 \label{eq:objective}
\end{equation}
The fixed divisors define the relative scale used to train the optional learned supervisor. The manual supervisor applies the timetable and budget laws directly; it does not minimize this composite by learning. The fast model minimizes predicted $W+R$. Actual passenger exposure, public dispatch delay, readiness, and CV$^2$ are evaluated separately under each study's frozen hypotheses.

At decision time $t$, the controller uses a history $\mathcal I_t$ consisting of available APC events, AVL telemetry, completed-trip summaries, and published plans. APC fields include time, station, direction, boarding and alighting counts, and entering load. Completion of a service event can occur after its arrival observation. Historical statistics requiring completed service are formed only after completion; currently available events can already inform queue reconciliation. Simulator passenger records are used to score physical outcomes, not to fill an unobserved current queue in the predictive graph.

The distinction between boarding and arrival rates is operationally consequential. A vehicle with spare capacity that serves the waiting passengers provides a clearance observation. A full vehicle gives a censored observation of demand: the number boarded does not determine the number remaining. Between clearance events, a model may propagate a queue using an estimated arrival rate, but that propagation remains a forecast. The controller cannot infer that no passenger arrived merely because no bus boarded passengers during the interval.

APC availability follows the simulator's recorded arrival-time boarding assignment. In contrast, statistics requiring service completion or a physical return enter their histories only at the corresponding completion clock. These two event clocks are retained in the executable observation model.

\section{Timetable and Holding-Budget Interface}
\label{sec:method}
\begin{figure*}[t]
\centering
\begin{tikzpicture}[font=\small,>=Latex,
 box/.style={draw=blue!55!black,rounded corners=2pt,fill=blue!4,align=center,minimum height=1.0cm,text width=3.0cm},
 arr/.style={->,thick,blue!55!black}]
\node[box] (data) at (0,0) {Available APC / AVL\\Closed history};
\node[box] (freq) at (4.1,0) {Causal harmonic state\\LF level / forecast\\Residual energy};
\node[box] (slow) at (8.2,0) {Slow policy: 900 s\\Budget or joint\\LF timetable law};
\node[box] (ledger) at (12.3,0) {Published releases\\Persistent trip budgets};
\node[box] (graph) at (4.1,-2.1) {Observed event graph\\APC service forecast\\Plan-aware continuation};
\node[box] (fast) at (8.2,-2.1) {Fast event controller\\Legal-bin enumeration\\Predicted $W+R$};
\node[box] (physical) at (12.3,-2.1) {One executed holding\\One budget debit\\Actual returns / departures};
\draw[arr] (data)--(freq);
\draw[arr] (freq)--(slow);
\draw[arr] (slow)--(ledger);
\draw[arr] (data.south) |- (graph.west);
\draw[arr] (ledger.south) -- ++(0,-.32) -| ([xshift=-.35cm]graph.north);
\draw[arr] (graph)--(fast);
\draw[arr] (fast)--(physical);
\draw[arr] (physical.south) -- ++(0,-.5) -| (data.south);
\end{tikzpicture}
\caption{Shared information and resource interface, reused from the development architecture. Manual joint fixes the slow mode to joint; the learned comparator selects budget or joint. The fast controller uses APC/AVL, published plans, and remaining budgets. Frequency summaries enter the slow state and timetable law. Physical events update observations and resource accounting.}
\label{fig:architecture}
\end{figure*}

\subsection{Causal Frequency Summaries}
The frequency tracker forms closed-bin APC signals and maintains a causal harmonic state. For a nonnegative observed rate $x_n$, let $z_n=\log(1+x_n)$ and let $\phi_n$ contain an intercept, a temporal trend, and harmonic basis terms. A recursive least-squares update has the form
\begin{align}
 k_n&=\frac{P_{n-1}\phi_n}{\rho+\phi_n^\top P_{n-1}\phi_n},\\
 \theta_n&=\theta_{n-1}+k_n(z_n-\phi_n^\top\theta_{n-1}),\\
 P_n&=\rho^{-1}(P_{n-1}-k_n\phi_n^\top P_{n-1}).
\end{align}
Here $\rho$ is the recursive estimator's forgetting factor, distinct from budget retention $\rho_B$. The transformed fitted value supplies a nonnegative low-frequency (LF) summary $L_n$. A residual is formed against the \emph{pre-update} prediction, and its smoothed value and high-frequency (HF) energy are retained. Thus the fitted low component and smoothed residual are not an exact additive reconstruction of the raw count. A future LF summary is obtained by evaluating the fitted state at the future harmonic phase, without incorporating future observations.

The implemented slow state includes global low demand, its slope, a low-frequency forecast, and high-frequency energy. The timetable rule uses the global low level and forecast. Station-specific short-window arrival-rate estimates in the fast graph are obtained separately from APC clearance intervals. They have units of passengers per second and are not interchangeable with boarding-bin residuals. This separation avoids adding quantities with different sampling and exposure meanings.

\subsection{Bounded Dynamic Timetable Revision}
The slow manager begins at 7200 s and considers publication every 900 s. For each direction, a revision concerns the next mutable departure and preserves the following departure as an anchor. Write $p_1$ for the next public release, $p_2$ for the following public release, and $d_-$ for the preceding actual departure. Let $L_t$ and $F_t$ be the currently supported low-frequency level and forecast. Define
\begin{align}
 \omega_t&=\clip((p_2-t)/1800,0,1),\\
 \ell_t&=\pos{L_t},\qquad r_t=\pos{(1-\omega_t)L_t+\omega_t F_t},\\
 \tau^*&=\frac{\ell_t d_-+r_t p_2}{\ell_t+r_t}.
 \label{eq:timetable}
\end{align}
If the frequency estimate is unsupported, equal weights replace $\ell_t$ and $r_t$. The proposed release is restricted to a $\pm30$ s shift around $p_1$, then constrained by the current time, a 180 s minimum separation from the preceding published release, and a 180 s separation from the anchor. The executable value is rounded upward to the simulator's integer-second clock. Both the revised trip's target headway and the anchor's following target headway are updated.

A revision requires two mutable future departures within the publication window, sufficient completed-trip history, and a compatible vehicle-supply estimate. A trip that has already become operationally committed is not rewritten. The rule evaluates supply from available vehicles and forecast returning vehicles; it does not create an additional vehicle or reserve an unobserved return. Manual joint always enables this law and the budget mechanism. The optional learned supervisor controls permission to invoke the law rather than directly optimizing timetable entries.

\subsection{Persistent Holding Budgets}
For each active vehicle--trip pair $i$, the manager estimates the normal remaining command exposure $C_i$ from the number of remaining controls and the median command seconds per control in at least three completed eligible trips of the same direction. A return estimate separates the remaining non-holding running and service component, irrevocable commands, and the physical overhead of control events. It is matched to an upcoming trip within a 2700 s supply look-ahead.

Let $\widehat r_i^{\mathrm{ord}}$ be the estimated ordinary return time, and $p_{j(i)}$ the public departure of the matched trip. The estimated supply pressure and proposed remaining allowance are
\begin{align}
 P_i&=\pos{\widehat r_i^{\mathrm{ord}}+1-p_{j(i)}},\\
 \widetilde B_i&=\left\lceil\max(\rho_B C_i,C_i-P_i)\right\rceil.
 \label{eq:budget}
\end{align}
The default uses $\rho_B=0.5$; the independent configuration study also evaluates 0.25 without retraining. This retention floor preserves a fraction of the estimated ordinary command exposure; a larger value permits more holding time at the same supply pressure. The ordinary exposure is estimated from the policy's own completed-trip history, not an externally fixed regularity requirement. The allowance is vehicle time, not a vehicle or driver fraction. Only active trips with positive pressure, positive normal exposure, and remaining control opportunities receive a new allowance. If one already exists, publication takes $\min(B_i,\widetilde B_i)$ rather than replenishing it. At final submission of a command,
\begin{equation}
 B_i\leftarrow B_i-h_e,\qquad h_e\leq B_i^{\mathrm{before}}.
 \label{eq:debit}
\end{equation}
The allowance is closed when that trip returns. It persists across changes in the upper mode and across all remaining control points.

Budget accounting uses command seconds. In the simulator, a zero command still occupies one physical tick at a control point. Thus physical residence is $\max(1,h_e)$ for a submitted command, with corresponding automatic zero ticks at unrequested controls. Treating residence as identical to the budget debit would change the executed model. The return predictor estimates this overhead separately; it does not claim that subtracting one budget second always advances a return by one second.

\subsection{Fast Holding with a Common Continuation}
The native bilateral request uses forward departure spacing $g_f$ and a causal estimate of following-bus action-ready spacing $g_b$. When both are supported, its continuous request is $0.25(g_b-g_f)/2$. Without following-bus evidence it uses $0.25(T-g_f)$ for the trip target $T$; without forward evidence it requests zero. The result is clipped to $[0,45]$ s and mapped to the nearest original legal bin, with ties selecting the shorter command. Follower identity and timing use current telemetry and completed observations, not future return labels.

At a controllable event, the fast layer constructs a graph from the same current observations, published plans, and budget ledger. Its reference is the native bilateral rule projected onto the remaining allowance. If the native command is $h_e^{\mathrm{nat}}$, the projection is
\begin{equation}
 \bar h_e=\max\{h\in\mathcal H_e:h\leq h_e^{\mathrm{nat}},\ h\leq B_i\},
 \label{eq:projection}
\end{equation}
where the budget restriction is omitted for an unbudgeted trip. The rule's original legal mask is preserved. The current controller enumerates the legal members of $\{0,5,10,15,20,30,45\}$ s, using the same future native continuation and a separate copy of the ledger for each evaluation. The selected action minimizes predicted $W+R$ over candidates that satisfy the frozen relative feasibility conditions: predicted CV$^2$, relevant vehicle return times, subsequent departure times, and trip departure delays must not exceed their reference values. Coverage must be complete for the required events. The first executable command is applied; future actions remain predictions.

Published releases affect terminal events and downstream supply matching in the graph. Budget depletion changes later available commands. These dependencies mean that the two layers interact through a shared feasible set, even if the immediate passenger-cost ordering of the current holding candidates does not change. Every trial starts from the same remaining allowance, so searching more candidates does not spend physical budget. Only the executed command debits the live ledger.

If the predictive reference lacks the required support, the controller uses the budget-feasible native command. This fallback is an operational part of the method. An enumeration can be exact over the current finite library and still fail to identify the best real-world action, because the future continuation and demand estimates are approximate. We therefore distinguish search completeness from prediction accuracy and from actual transport benefit.

\subsection{Optional Learned Publication-Permission Comparator}
The upper action set is $\mathcal A=\{\mathrm{budget},\mathrm{joint}\}$. Both modes permit persistent budgets; only the joint mode permits a new timetable revision. The 180-dimensional state contains the time, fleet load, frequency summaries, accumulated arrival-headway moments, public-delay debt, the previous mode, directional supply and plan summaries, completed-trip history, and telemetry and remaining allowances for the 12 vehicles. Fixed divisors normalize these fields. The same canonical action labels are used in the actor, critic, replay, previous-action encoding, and execution.

The reward is the negative increment of~\eqref{eq:objective} over a slow interval, with $\gamma=1$ and the final interval closed after the service tail. A discrete maximum-entropy actor--critic uses twin critics, a 64-unit hidden width and fixed entropy temperature $10^{-5}$. Evaluation is deterministic. The complete optimization settings and update equations are provided in the Supporting Information; this optimizer is a matched comparator, not the proposed contribution.

\begin{algorithm}[t]
\caption{The shared timetable--holding interface}
\label{alg:control}
\begin{algorithmic}[1]
\STATE Update histories from newly available APC/AVL events.
\IF{a new slow publication time has been reached}
\STATE Use joint for manual control, or select budget/joint with the frozen supervisor.
\STATE Issue or tighten trip budgets using~\eqref{eq:budget}.
\IF{joint is selected and a plan revision is supported}
\STATE Publish the bounded release and its two target headways.
\ENDIF
\ENDIF
\IF{a vehicle requests a holding action}
\STATE Build the observed graph with current plans and budgets.
\STATE Project the native reference using~\eqref{eq:projection}.
\STATE Enumerate current legal candidates with separate ledgers.
\STATE Select the lowest-cost feasible action or reference fallback.
\STATE Submit one command and debit its live trip allowance once.
\ENDIF
\STATE Close the allowance when its physical trip returns.
\end{algorithmic}
\end{algorithm}

\subsection{Joint Credit Admission as a Development Candidate}
Pressure credit in~\eqref{eq:budget} restricts the continuation before the fast model checks relative feasibility. We therefore evaluate a separate candidate that treats new credit itself as a decision. At the slow clock, the original formula proposes an allowance without imposing it. At the proposing cycle's next supported holding query, the unchanged fast selector with existing credit produces a common no-new-tightening reference. We enumerate current legal commands under the proposed credit and require their continuations to meet the same reference's passenger-cost, CV$^2$, return, next-launch, and dispatch-delay conditions. The lowest predicted passenger cost is preferred; exact cost ties prefer lower total public delay. Only a strict cost or delay improvement admits the new allowance.

Commands submitted between proposal and admission debit the proposed remainder. A request expires at the next publication or physical return. Unsupported evidence imposes no new tightening; existing credit and irreversible commands remain unchanged. The chosen allowance is committed before its current command is submitted. This candidate changes credit admission, not the demand predictor, legal bins, cost weights, or fleet.

\subsection{Operational Coupling}
The architecture makes two interactions explicit: slow releases alter fast resource feasibility, and fast holding consumes vehicle time that constrains later releases. It also exposes a useful distinction between \emph{permission} and \emph{dose}. Selecting joint does not guarantee that a plan will be revised; the law may reject a publication, produce the same integer-second release, or be unable to influence an already committed trip. An upper action count is therefore insufficient evidence of timetable control. We record published revisions, binding to actual departures, and changes relative to matched control paths.

The implemented frequency pathway runs from the causal summary to the slow state or timetable law, and then from published plans and budgets to fast control. The architecture in Fig.~\ref{fig:architecture} depicts these paths. The representation and supervisor comparisons evaluate their incremental effects separately from the budget-setting contrast.
